\documentclass[11pt,a4paper]{article}
\usepackage[T1]{fontenc}
\usepackage[utf8]{inputenc}
\usepackage{amsmath,amssymb,amsthm}
\usepackage[margin=1in]{geometry}
\usepackage[colorlinks=true,linkcolor=blue,urlcolor=blue]{hyperref}
\theoremstyle{plain}
\newtheorem{theorem}{Theorem}
\newtheorem{lemma}[theorem]{Lemma}
\newtheorem{proposition}[theorem]{Proposition}
\newtheorem{corollary}[theorem]{Corollary}
\theoremstyle{definition}
\newtheorem{remark}[theorem]{Remark}

\title{Recovering the unwritten recurrences of the table OEIS A237945}
\author{Adrian Perez Fontelles\\ \small Independent researcher}
\date{2 September 2026}
\begin{document}
\maketitle

\begin{abstract}
OEIS A237945 is a two-dimensional table $T(n,k)$ whose empirical recurrences are, for several of
its lines, given only as an ORDER --- ``k=2: [order 12]'' and the like --- with the
coefficients in a linked file rather than in the entry. Those conjectures can still be settled,
and the recurrences recovered. A line of the table is a fixed-width or fixed-height array
count, hence a walk count on finitely many vertices, hence satisfies some monic recurrence of
order at most that number; Berlekamp--Massey on twice as many exact terms returns the MINIMAL
one. Where its order equals the order the entry states --- and where the same holds after the
first few terms are dropped --- any recurrence of that order the line satisfies has a
characteristic polynomial that is a multiple of the minimal one and of equal degree, hence is
that polynomial. This note settles column 2, column 3, column 4.
\end{abstract}

\noindent\small 2020 Mathematics Subject Classification. 05A15, 11B37, 15A18.\normalsize

\section{The table and the unwritten conjectures}

OEIS A237945 is ``T(n,k)=Number of (n+1)X(k+1) 0..2 arrays with the maximum minus the minimum of every 2X2 subblock equal.''. It is read by antidiagonals and begins
\[
81, 437, 437, 2417, 5519, 2417, 13893, 70435,\ \dots
\]
The entry carries, among its empirical claims:
\begin{quote}\small
k=2: [order 12]\\
k=3: [order 23]\\
k=4: [order 54]
\end{quote}
As of the ``Last modified'' line on the live entry (Jul 23 2025 10:37:35, revision 6) these are still
recorded as empirical, and the recurrences themselves are not in the entry text.

\section{Why a line of the table is a walk count}

Fix the index that the line fixes. The entry defines $T(n,k)$ as a number of arrays of a shape
determined by $n$ and $k$, subject to a condition local to a bounded window of consecutive
lines. Substituting the fixed index into the entry's own wording turns the definition into an
ordinary one-parameter array count, and for such a count the admissible windows are the
vertices of a finite digraph and an array is a walk. So the line is a walk count on $S$
vertices, and by the Cayley--Hamilton theorem it satisfies a monic integer recurrence of order
at most $S$.

\section{Recovering the recurrences}

\begin{lemma}\label{lem:bm}
A sequence known to satisfy some linear recurrence of order at most $S$ is determined, with its
minimal recurrence, by its first $2S$ terms; Berlekamp--Massey applied to them returns that
minimal recurrence exactly.
\end{lemma}

\subsection*{Column $k=2$}
Substituting the fixed index into the entry's wording gives an ordinary one-parameter array
count, a walk on $S=47$ vertices. Berlekamp--Massey on $2S=94$ exact terms
returns a minimal recurrence of order $12$ --- the order the entry states --- with
integer coefficients
\begin{center}\small\begin{tabular}{cccccccc}
$c_{1}$ & $20$ & $c_{4}$ & $1719$ & $c_{7}$ & $8016$ & $c_{10}$ & $1560$ \\
$c_{2}$ & $-56$ & $c_{5}$ & $2496$ & $c_{8}$ & $10172$ & $c_{11}$ & $1344$ \\
$c_{3}$ & $-524$ & $c_{6}$ & $-12490$ & $c_{9}$ & $-11936$ & $c_{12}$ & $-320$ \\
\end{tabular}\end{center}
so that $a(m)=\sum_{i=1}^{12}c_i\,a(m-i)$ for every $m$. Removing the first
$12+4$ terms and repeating gives order $12$ again, which is what the
identification below needs.

\subsection*{Column $k=3$}
Substituting the fixed index into the entry's wording gives an ordinary one-parameter array
count, a walk on $S=125$ vertices. Berlekamp--Massey on $2S=250$ exact terms
returns a minimal recurrence of order $23$ --- the order the entry states --- with
integer coefficients
\begin{center}\small\begin{tabular}{cccccccc}
$c_{1}$ & $39$ & $c_{7}$ & $-11041734$ & $c_{13}$ & $1905197584$ & $c_{19}$ & $-377080192$ \\
$c_{2}$ & $-58$ & $c_{8}$ & $26665646$ & $c_{14}$ & $-4027531552$ & $c_{20}$ & $1576797696$ \\
$c_{3}$ & $-7030$ & $c_{9}$ & $125943418$ & $c_{15}$ & $-2766833480$ & $c_{21}$ & $-161140224$ \\
$c_{4}$ & $21065$ & $c_{10}$ & $-291658456$ & $c_{16}$ & $5675870464$ & $c_{22}$ & $-210014208$ \\
$c_{5}$ & $435029$ & $c_{11}$ & $-685365122$ & $c_{17}$ & $1922313440$ & $c_{23}$ & $47554560$ \\
$c_{6}$ & $-1133100$ & $c_{12}$ & $1522606472$ & $c_{18}$ & $-4271600256$ &  &  \\
\end{tabular}\end{center}
so that $a(m)=\sum_{i=1}^{23}c_i\,a(m-i)$ for every $m$. Removing the first
$23+4$ terms and repeating gives order $23$ again, which is what the
identification below needs.

\subsection*{Column $k=4$}
Substituting the fixed index into the entry's wording gives an ordinary one-parameter array
count, a walk on $S=345$ vertices. Berlekamp--Massey on $2S=690$ exact terms
returns a minimal recurrence of order $54$ --- the order the entry states --- with
integer coefficients
\[
(c_1,\dots,c_{54})=(114,\ -2867,\ -52231,\ 2711827,\ -8661804,\ -757325894,\ 7527306920,\ 79807056462,\ -1367998363447,\ -1687663009317,\ 110917108096881,\ -258739310124081,\ -4739655400616131,\ 22152458780958445,\ 110002571126724549,\ -829556868674006798,\ -1130920306950853929,\ 18146530385042371408,\ -6950167137661256282,\ -250068856733869338602,\ 386266227835423261411,\ 2190664090388691523297,\ -5761584646668520683269,\ -11442161936878462930223,\ 49132704115136727028958,\ 24394672643387282353612,\ -266910375387161278937832,\ 97283434548590124500248,\ 938212562586266131806940,\ -959939086000441677257788,\ -2038816696572809784140896,\ 3609791640488783037461504,\ 2220416381017563009980720,\ -7837086690644302787241200,\ 610276958764461994800416,\ 10303905957021762786240576,\ -5548883776251089714762048,\ -7713746166965585532004736,\ 7880574088008778526416384,\ 2354164560770143102939648,\ -5548434853033827411120128,\ 757783782272487607109632,\ 2016336047968048902709248,\ -855816501085111798521856,\ -303435904087044273700864,\ 257929269161045172879360,\ -8404836223184201383936,\ -30686317992187367522304,\ 5967532860739167715328,\ 1322995206744814125056,\ -477429732174235435008,\ 3181318536299544576,\ 11381934187654152192,\ -994740340512522240).
\]
so that $a(m)=\sum_{i=1}^{54}c_i\,a(m-i)$ for every $m$. Removing the first
$54+4$ terms and repeating gives order $54$ again, which is what the
identification below needs.

\begin{theorem}
For each line above, the displayed recurrence holds for every index, and it is the recurrence
the entry records.
\end{theorem}

\begin{proof}
It holds by Lemma~\ref{lem:bm}. For the identification, the entry asserts that the line
satisfies SOME recurrence of the stated order, possibly only from some index on. Let $q$ be its
characteristic polynomial and $q_{\min}$ the minimal polynomial of the tail. Then
$q_{\min}\mid q$, and the second Berlekamp--Massey run --- on the line with its first few
terms removed --- shows $\deg q_{\min}$ equals the stated order, which is $\deg q$. Both are
monic, so $q=q_{\min}$.
\end{proof}

\section{Verification}

Three checks.

First, each line's model was compared against the entry's own published values for that line,
taken from the antidiagonal DATA read in either orientation and from the ``Table starts''
block, and agrees at every available term. This is the only point at which reading the entry's
English enters, and a misreading gives different counts.

Second, each recovered recurrence was evaluated directly on the model's values, with no
Berlekamp--Massey involved, and holds at every index tested.

Third, each order was computed twice by independent routes --- modulo a $61$-bit prime, where
every intermediate is a machine word, and again in exact rational arithmetic --- and once more
on the line with its first few terms dropped, which is what the identification requires. All
agree.

\begin{thebibliography}{9}
\bibitem{oeis} The OEIS Foundation, \emph{The On-Line Encyclopedia of Integer
Sequences}, \url{https://oeis.org}, 2026. Sequence A237945.
\bibitem{massey} J.~L.~Massey, Shift-register synthesis and BCH decoding, \emph{IEEE
Trans. Inform. Theory} \textbf{15} (1969), 122--127.
\bibitem{stanley} R.~P.~Stanley, \emph{Enumerative Combinatorics}, Volume 1, 2nd ed.,
Cambridge University Press, 2011. (Transfer-matrix method, Section 4.7.)
\bibitem{horn} R.~A.~Horn and C.~R.~Johnson, \emph{Matrix Analysis}, 2nd ed., Cambridge
University Press, 2013.
\end{thebibliography}

\end{document}
